% =====================================================================================================
\section{How to read this document}\label{sec:intro}

This edition of the manuscript is written for one reader with two questions: \emph{do the assumptions make sense?} and
\emph{do the theorems look natural and intelligible?} It follows the method from the data to the guarantees, in this
order.
\begin{enumerate}[leftmargin=1.6em,itemsep=1pt]
  \item \emph{The model and the priors} (Section~\ref{sec:method}): what the software assumes and sets, every default with
    the reason for it (Table~\ref{tab:prior}).
  \item \emph{The sampler} (Section~\ref{sec:sampler}): the algorithm as coded, how to read its outputs, which parts the
    theorems cover (Table~\ref{tab:coverage}), and, separately, the parallel engine built for $p=10^6$, which no theorem
    covers.
  \item \emph{The assumptions} (Section~\ref{sec:assumptions}): one table of verdicts (Table~\ref{tab:a-summary}), then one
    entry per assumption with five parts: (a) the statement, (b) its meaning in words, (c) what it rules out, (d) whether it
    is reasonable for real designs, and (e) how it compares with the conditions of \citet{SongLiang2015}, the split-and-merge
    paper.
  \item \emph{The proved results} (Section~\ref{sec:proved}), the main ones first. Each comes in words, then as a formal
    statement, then with a grey status box quoting the project's registry, a paragraph on why it is natural, and a
    comparison with \citet{SongLiang2015}.
  \item \emph{The planned results} (Section~\ref{sec:planned}), each marked \textsf{PLANNED}: nothing there is claimed.
  \item \emph{The statuses} (Section~\ref{sec:statuses}): what changed in the registry since the theory freeze of
    2026-11-17, and the open questions.
\end{enumerate}
The supplement (Appendix~\ref{sec:supp}) holds the minor results, listed first with one line each, and the appendices hold
the notation, the full parameter table with code locations, the conventions in full and the status ledger. The 1M runs and
the benchmarks are reported elsewhere.

\paragraph{The method in one paragraph.} SMC-BVS puts a fixed prior on models and coefficients: a beta-binomial prior on
the model size \citep{ScottBerger2010} and a $g$-prior with a unit ridge, so that the posterior over models is known in
closed form up to its normalising constant, the evidence. Nothing is tuned by the user. The method partitions the
predictors into small blocks of correlated predictors, scores every subset of every block exactly, and combines these
block tables into a \emph{start law}, an approximation of the posterior that can be computed and sampled exactly. A
sequential Monte Carlo sampler, run as four independent \emph{islands}, then moves from the start law to the posterior
along a tempered path that an internal calibration fixes in advance. Two deterministic local searches, one of them from
the mode of the start law, finish the run. The run reports its best model, pooled inclusion probabilities, their median
probability model and an evidence estimate.

\paragraph{What is proved, in one paragraph.} Under the assumptions of Section~\ref{sec:assumptions}: the posterior puts
all but a fraction of order $\rn$ of its mass on the true model (Result~1), and recovers the true conditional law of the
response at that radius, with no lower bound on the coefficients (Result~2). For block and AR designs whose true
coefficients do not cancel, every run that produces output reports exactly the true model, with probability tending to
one, for every seed (Result~3). Without those design conditions, a run rarely certifies a small model that misses a true
predictor (Result~4). The evidence estimate is unbiased on every design (Result~5). The rate $\rn$ cannot be improved
(Result~6). All results are asymptotic in $n$, with $p$ at most polynomial in $n$; this paper contains no simulation study.

\paragraph{What is not claimed.} No theorem covers the parallel engine used at $p=10^6$, more than four islands, or a run
that stops early. Statements about a run's estimates concern the event that its calibration succeeds, or that the run
completes; no lower bound on the probability of that event is claimed, except that under Convention~\as{C2p} every run
produces output. The consistency of the pooled inclusion probabilities and evidence estimate is proved only on designs
without leakage (supplement, Remark~\ref{thm:outputs}).
